\documentclass[11pt]{article}

\usepackage{nmrdoc}
\usepackage{siunitx}
\usepackage{bm}

\sisetup{per-mode=symbol}
\DeclareSIUnit{\angstrom}{\text{\AA}}
\DeclareSIUnit{\molar}{M}
\setlength{\parindent}{1.2em}
\emergencystretch=2em

\begin{document}
\NmrTitle{The APBS Electric-Field Calculator}

\section{What the calculator computes}

The ApbsField calculator provides the electrostatic quantities used for the
Buckingham electric-field contribution to nuclear shielding. Buckingham's idea
is that an electric field at a nucleus can perturb the local electronic
response, and the perturbed electrons then change the magnetic shielding felt by
that nucleus. This implementation stops at the field-derived descriptor. It
does not multiply by Buckingham response constants and it does not emit a
shielding in ppm.

For each atom \(i\), the code stores a field vector and a field-gradient
matrix. The first is the solvated electric field
\begin{equation}
  \bm E_i = (E_x,E_y,E_z)
\end{equation}
in \(\si{\volt\per\angstrom}\). The second is a rank-2 electric-field-gradient
matrix
\begin{equation}
  G_{ab,i} = \frac{\partial E_a}{\partial r_b}(\bm r_i),
  \qquad a,b\in\{x,y,z\},
  \label{eq:g-def}
\end{equation}
in \(\si{\volt\per\angstrom\squared}\). The row index \(a\) names the electric
field component and the column index \(b\) names the coordinate direction in
which the field is differentiated. The code symmetrizes this matrix and removes
its trace before the spherical-tensor split, so the feature exported for
\(\bm G_i\) is the five-component traceless symmetric part. In this document,
``kernel'' means this uncalibrated field-derived tensor quantity: it is the
geometric and electrostatic input that can correlate with a DFT shielding
contribution after calibration.

\section{Where shielding enters}

For a nucleus in an applied magnetic field \(\bm B_0\), the local magnetic field
is written
\begin{equation}
  \bm B_{\mathrm{nuc}} = (\bm I-\bm\sigma)\bm B_0 .
  \label{eq:shielding}
\end{equation}
\(\bm I\) is the \(3\times3\) identity matrix. \(\bm\sigma\) is the nuclear
shielding tensor, a dimensionless linear-response coefficient: it tells how the
electrons reduce or redirect the applied field at the nucleus. A matrix is
needed because a molecule has orientation. A field applied along one laboratory
axis can induce an electronic response whose field at the nucleus has
components along several axes.

The electronic response in Eq.~\eqref{eq:shielding} is quantum mechanical. The
calculator does not solve that response. Instead it supplies a classical
electrostatic descriptor for one route by which the response changes. In an
inhomogeneous protein environment, nearby charges, dielectric screening, and
ions create an electric field at the nucleus. Buckingham-type expansions model
the shielding perturbation as depending on that field and on how the field
varies across the electronic cloud. A common schematic form is
\begin{equation}
  \Delta\sigma_{\mathrm{iso}}
    \approx A_X E_\parallel + B_X E_\parallel^2,
  \qquad
  \Delta\bm\sigma_{2}
    \approx \gamma_X\,\bm G_{2}.
  \label{eq:buckingham-sketch}
\end{equation}
\(X\) labels the nucleus or atom type; \(A_X\), \(B_X\), and \(\gamma_X\) are
calibrated response coefficients carrying the units needed to convert
field units into shielding units; \(E_\parallel\), in
\(\si{\volt\per\angstrom}\), is a chosen projection of the electric field, for
example along a local bond axis; and \(\bm G_2\), in
\(\si{\volt\per\angstrom\squared}\), is the traceless symmetric rank-2 part of
the field gradient. Equation
\eqref{eq:buckingham-sketch} explains why ApbsField keeps both a vector field
and a rank-2 tensor. It is not the formula evaluated by this C++ calculator.

The APBS part matters because the electrostatic field is not the vacuum Coulomb
sum over protein partial charges. The APBS solve includes a dielectric boundary
and mobile-ion screening. ApbsField reads the resulting potential grid through
the C bridge in \texttt{src/apbs\_bridge.h}; the internal numerical method used
by APBS is outside this calculator's model.

\section{The tensor split used by the code}

The shielding tensor in Eq.~\eqref{eq:shielding} and the field-gradient tensor
in Eq.~\eqref{eq:g-def} are both rank-2 objects: each maps one direction index
to another. The project stores such a \(3\times3\) tensor \(\bm X\) by a
closed-form split into scalar, antisymmetric, and symmetric-traceless parts. The
C++ owner of this convention is \texttt{SphericalTensor::Decompose} in
\texttt{src/Types.cpp}.

The scalar part is
\begin{equation}
  T_0 = \frac{1}{3}\operatorname{tr}\bm X .
  \label{eq:t0}
\end{equation}
It is the isotropic part: the average of the three diagonal entries. For a
shielding tensor this is the part that survives rapid molecular tumbling and
sets the isotropic shielding used in ordinary solution shifts after reference
subtraction.

The antisymmetric part is stored as a three-component pseudovector,
\begin{equation}
  \bm T_1 =
  \frac{1}{2}
  \begin{pmatrix}
    X_{yz}-X_{zy}\\
    X_{zx}-X_{xz}\\
    X_{xy}-X_{yx}
  \end{pmatrix}.
  \label{eq:t1}
\end{equation}
This is the transpose-odd part of the matrix written in vector form. It is
usually not observed directly in standard isotropic NMR, but the representation
keeps it when a calculator produces it.

The symmetric traceless matrix is
\begin{equation}
  \bm S =
  \frac{\bm X+\bm X^{\mathsf T}}{2}
  - T_0\bm I ,
  \qquad \operatorname{tr}\bm S=0 .
  \label{eq:sym-traceless}
\end{equation}
It has five independent components. The code packs those components into the
real spherical \(T_2\) basis
\begin{equation}
  \bm T_2 =
  \begin{pmatrix}
    \sqrt{2}\,S_{xy}\\
    \sqrt{2}\,S_{yz}\\
    \sqrt{3/2}\,S_{zz}\\
    \sqrt{2}\,S_{xz}\\
    (S_{xx}-S_{yy})/\sqrt{2}
  \end{pmatrix}.
  \label{eq:t2}
\end{equation}
The numerical factors in Eq.~\eqref{eq:t2} are the code's normalization. They
make the squared length of the five-vector equal to the Frobenius norm
\(\sum_{ab}S_{ab}^2\) of the underlying matrix. The full nine-double layout is
\begin{equation}
  [T_0,T_{1x},T_{1y},T_{1z},
    T_{2,-2},T_{2,-1},T_{2,0},T_{2,+1},T_{2,+2}] .
\end{equation}
ApbsField stores the Cartesian field-gradient matrix on the atom, decomposes it
with this routine, and writes the \(T_2\) five-vector. Because the matrix has
already been symmetrized and de-traced, \(T_0\) and \(T_1\) are structural
zeros in the exported APBS field-gradient feature.

\section{The method implemented}

ApbsField depends on \texttt{ChargeAssignmentResult}. That upstream result
places a partial charge \(q_j\), in elementary charges, and a Poisson--Boltzmann
radius \(a_j\), in \(\si{\angstrom}\), on each atom from the active charge
source. In the ff14SB flat-table path those values come from
\texttt{data/ff14sb\_params.dat}; AMBER topology and preloaded-charge paths
can also provide the same fields. If a radius is missing, the APBS calculation
returns failure. If the charge table reports placeholder PB radii, the code
still fills the APBS fields but logs that the dielectric boundary is not
physically validated.

The calculator forms separate coordinate, charge, and radius arrays and passes
them to \texttt{apbs\_solve}. The configured values in
\texttt{data/calculator\_params.toml} are
\begin{equation}
\begin{aligned}
  N_{\mathrm{grid}} &= 161 \quad \text{points per axis},\\
  F_d &= \max(L_d+\SI{40.0}{\angstrom},\,\SI{40.0}{\angstrom}),\\
  C_d &= F_d+\SI{30.0}{\angstrom},\\
  \epsilon_{\mathrm{protein}} &= 4.0,\qquad
  \epsilon_{\mathrm{solvent}} = 78.54,\\
  T &= \SI{298.15}{\kelvin},\qquad
  I = \SI{0.15}{\molar}.
\end{aligned}
\label{eq:apbs-params}
\end{equation}
\(L_d\) is the protein bounding-box extent along Cartesian direction \(d\).
The C++ layer computes both fine-grid lengths \(F_d\) and coarse-grid lengths
\(C_d\). The APBS bridge used here performs a single linearized
Poisson--Boltzmann solve on the coarse box \(C_d\), with single-Debye--Huckel
boundary conditions, and returns a scalar potential grid \(\phi\) in units of
\(kT/e\). The bridge also fixes the mobile ions as a \(+1/-1\) pair at
concentration \(I\), with radii \(\SI{0.95}{\angstrom}\) and
\(\SI{1.81}{\angstrom}\), and passes a solvent probe radius of
\(\SI{1.4}{\angstrom}\) to APBS. It does not return an electric field or a
field gradient.

At an atom position \(\bm r_i\), ApbsField first interpolates \(\phi\) from the
grid by trilinear interpolation, the eight-corner weighted average inside the
cell containing the requested point. If the requested cell falls outside the
returned grid, the interpolation helper returns \(0.0\); the padding values are
intended to keep atom stencils inside the grid. Let \(h_a\) be the APBS grid
spacing along axis \(a\), and let \(\hat{\bm e}_a\) be the unit vector in that
direction. The electric field is then computed by a centered difference,
\begin{equation}
  E_a(\bm r_i)
  =
  -\,\frac{
    \phi(\bm r_i+h_a\hat{\bm e}_a)
    -
    \phi(\bm r_i-h_a\hat{\bm e}_a)}
  {2h_a}.
  \label{eq:central-e}
\end{equation}
The minus sign is the electrostatic convention \(\bm E=-\nabla\phi\). Before
unit conversion, Eq.~\eqref{eq:central-e} has units \(kT/(e\,\si{\angstrom})\).

The field-gradient matrix is another centered difference:
\begin{equation}
  G_{ab}(\bm r_i)
  =
  \frac{
    E_a(\bm r_i+h_b\hat{\bm e}_b)
    -
    E_a(\bm r_i-h_b\hat{\bm e}_b)}
  {2h_b}.
  \label{eq:central-g}
\end{equation}
This quantity has units \(kT/(e\,\si{\angstrom\squared})\) before conversion.
Since \(\bm E=-\nabla\phi\), \(\bm G\) is the negative Hessian of the APBS
potential. Some electrostatics texts use the opposite-sign Hessian as the
electric field gradient; this document follows the code's
\(\partial E_a/\partial r_b\) convention.
Equations \eqref{eq:central-e} and \eqref{eq:central-g} define the local
finite-difference model the calculator keeps from the APBS potential. Near atom
\(i\), after symmetrization, the potential has the Taylor expansion
\begin{equation}
  \phi(\bm r_i+\bm s)
  =
  \phi(\bm r_i)
  - \bm E_i\cdot\bm s
  - \frac{1}{2}\bm s^{\mathsf T}\bm G_i\bm s
  + O(|\bm s|^3).
  \label{eq:taylor}
\end{equation}
The field vector is the local slope of the potential with the electrostatic
minus sign; the matrix \(\bm G_i\) records how that slope changes for small
displacements. This quadratic patch summarizes the computation:
ApbsField discards the absolute potential value and keeps the first- and
second-derivative information that Buckingham-type models can use.

The raw finite-difference matrix is then projected in two steps,
\begin{equation}
  \bm G^{\mathrm{sym}} =
  \frac{\bm G+\bm G^{\mathsf T}}{2},
  \qquad
  \widetilde{\bm G} =
  \bm G^{\mathrm{sym}}
  - \frac{\operatorname{tr}\bm G^{\mathrm{sym}}}{3}\bm I .
  \label{eq:project}
\end{equation}
Symmetrization removes antisymmetric residue from interpolation and differencing;
without it, the spherical decomposition would report a numerical \(T_1\)
pseudovector. The trace is the divergence of the electric field. For external
sources in a charge-free neighborhood, Laplace's equation gives zero trace. The
APBS grid also contains each atom's own Coulomb potential, whose point-source
Laplacian is represented on the grid as a large finite trace near the atom.
Subtracting \(\operatorname{tr}\bm G/3\) from the diagonal removes that
self-potential trace artifact and leaves the traceless tensor exported by the
calculator.

Several guards are applied before storage. If the electric field contains a
NaN or infinity, both \(\bm E\) and \(\widetilde{\bm G}\) are set to zero. If
the APBS-native field magnitude exceeds \(100.0\,kT/(e\,\si{\angstrom})\), the
field vector is rescaled to that magnitude; the field-gradient matrix is not
rescaled by that guard. Non-finite entries of the field-gradient matrix are set
to zero individually. Finally, both quantities are multiplied by the fixed
thermal voltage
\begin{equation}
  \frac{kT}{e} = \SI{0.025693}{\volt}
\end{equation}
from \texttt{src/PhysicalConstants.h}, matching the configured
\(\SI{298.15}{\kelvin}\), so the stored units become
\(\si{\volt\per\angstrom}\) for \(\bm E\) and
\(\si{\volt\per\angstrom\squared}\) for \(\widetilde{\bm G}\).

The feature export follows the stored quantities. For \(N\) atoms,
\texttt{apbs\_E.npy} has shape \((N,3)\) and stores the Cartesian electric
field. \texttt{apbs\_efg.npy} has shape \((N,5)\) and stores the five
\(T_2\) components of \(\widetilde{\bm G}\). If the APBS solve fails,
\texttt{Compute} returns
\texttt{nullptr}; the code does not substitute a vacuum Coulomb field.

\section{External inputs and output}

Two external ingredients feed this calculator. \texttt{ChargeAssignmentResult}
provides per-atom partial charges and PB radii from the active force-field
charge source, such as the ff14SB table or an AMBER topology; that source is
outside the ApbsField calculation. APBS takes those arrays and solves the
linearized Poisson--Boltzmann problem for the solvated potential; its solver
internals are behind \texttt{src/apbs\_bridge.h}. ApbsField consumes the
returned potential grid, not APBS-internal fields or gradients, and derives
\(\bm E\) and \(\widetilde{\bm G}\) itself. MOPAC charges, AIMNet2 charges or
embeddings, and DSSP secondary-structure annotations do not enter this
calculator.

The output is therefore an external-field descriptor: a solvated electric field
and a traceless symmetric field-gradient kernel at each atom. It bears on
shielding because Buckingham-type response terms can use the field vector and
the \(T_2\) field gradient as calibrated inputs. Before that calibration step,
the numbers are in \(\si{\volt\per\angstrom}\) and
\(\si{\volt\per\angstrom\squared}\), not ppm, and they should be read as
field-derived geometry rather than as predicted chemical shifts.

\end{document}
